\newpage
% \color{white}
% \pagecolor{black}

\ifdefined\useChapters
\chapter{Sempre aqui}
\else
\chapter{}
\fi

Quando fui para o futuro, vi muita merda acontecendo, muita gente usando suas habilidades e destruindo tudo ao seu redor. De tudo isso, o que mais me marcou foi ver nos rostos deles raiva e ódio. Ao contrário do que eu imaginava, com as habilidades, as pessoas não estavam felizes.

A felicidade é um dos termômetros de que estamos fazendo o que é certo, e isso é o que me motiva a não me esconder, mas a lutar contra esse futuro catastrófico, aconteça o que acontecer comigo.

Não quero juntar habilidades como quem junta dinheiro. Não preciso ser o mais forte, o mais temido, o mais raro. Até porque, usar todas as habilidades deveria ser um direito de todos. Imagina como o mundo seria incrível se cada um pudesse usar suas habilidades sem precisar humilhar ou se aproveitar dos outros.  
Mas alguém já teve essa ideia com os bens que temos hoje em dia — e, ainda assim, o resultado foi desigual. Deve haver um motivo. Por ora, preciso aprender o máximo que eu puder.

Já encontrei pessoas que movem objetos à distância. Por um tempo, também consegui. Conheci a menina que solta fogo pelas mãos, a que cura, o cara que volita, o que faz os outros perderem suas habilidades.  
E senti uma que me assustou de verdade: a que invade a cabeça.  
De todas essas, percebi uma coisa — ninguém tem duas habilidades, e todos passaram por algo triste antes de despertar a sua.  
Talvez o poder venha sempre de uma dor.

As habilidades são como sentidos, mas algumas são devastadoras. Outras, mais sutis. Todas valiosas — principalmente agora, quando preciso estar pronto para o próximo ataque.

Não sei por qual caminho começar. Nem quanto tempo tenho. Aquele futuro que vi pode ser daqui a dez anos... ou semana que vem.  
Não consegui reconhecer rostos, mas o desespero era real.

Mesmo que eu aprenda a controlar várias habilidades, ainda existe um problema: sozinho, não sei se consigo enfrentar todos os que seguem Crispim.  
Quem poderia me ajudar?  
Preciso de alguém forte. Ou, pelo menos, alguém em quem eu possa confiar.

A velha da floresta.  
Ela é a única que já me ajudou, e talvez seja a única que possa me guiar agora.

Enquanto pensava em como entrar novamente na floresta e alcançar o vilarejo, uma sensação estranha tomou minha cabeça — a mesma de quando espionei aquela festa, mas em menor intensidade.  
O mundo girou devagar, e senti o coração acelerar. Era aquele desespero familiar, o pânico do desconhecido, como se alguém encostasse dentro da minha mente.  
Só que, dessa vez, não havia dor. Havia calma.

Depois de um tempo, percebi: era como se uma voz falasse dentro da minha cabeça.

— Eu estou sempre aqui — disse uma voz feminina, suave, próxima demais.

Meu corpo gelou. Olhei em volta, procurando a origem do som, mas estava sozinho.

— Quem é você? Como está fazendo isso? — perguntei, quase sussurrando.

— Eu já te disse que você duvida demais, mesmo depois de ver e passar por tanta coisa — respondeu a voz, com um tom doce, quase divertido.

Aquela entonação me era familiar. Um arrepio percorreu o braço.

— É você? A senhora do vilarejo?

— Então era você aquele dia... — ela respondeu. — Você já vem me encontrando há muito tempo, mesmo sem perceber.

— Da outra vez, não tive tempo de perguntar o seu nome.

— Meu nome não faz diferença, filho. — A voz soava como se sorrisse. — E sim, venho te acompanhando há mais tempo do que imagina.

Por um instante, pensei ter sentido o ar vibrar, como se ela respirasse comigo.

— Já parou pra pensar que, às vezes, quando você tem uma ideia que não parece sua, talvez ela realmente não seja? — completou.

Fiquei em silêncio, tentando entender.  
— Quer dizer que você estava na minha cabeça esse tempo todo?

— Nem sempre sozinha. — O tom dela ficou mais grave, distante. — É certo que eu sempre disputei te influenciar com outras vozes. Mas isso não vem ao caso agora.

Aquela confissão me arrepiou. Eu não estava completamente sozinho, mas também não sabia quem mais me escutava.

— Então você já sabe o que está acontecendo e que eu preciso da sua ajuda.

— Eu sei — disse ela com calma. A serenidade era tão profunda que parecia me abraçar por dentro.  
— Mas será que você realmente entendeu o que está acontecendo? Você sabe o que está procurando?

— Eu... acho que sim. Tenho que impedir aquele futuro. Tenho que lutar.

— Lutar — repetiu ela, como se provasse o gosto da palavra. — Isso não é uma batalha de poderes, dessas que você vê nos filmes. E vencer não significa derrubar ninguém. Se fizer isso, estará apenas repetindo a verdade deles.

— Então o que eu tenho que fazer? — perguntei, irritado. — Você está me confundindo mais do que ajudando.

— Eu sei que esperava um parceiro de batalha, mas te digo: a luta que importa é travada na consciência. E só será vencida — se é que pode ser vencida — quando você entender e aceitar a natureza de tudo.

A voz foi diminuindo até virar um eco. O ar pareceu se expandir dentro da minha cabeça.  
Por um momento, pensei ter ouvido um último sussurro, quase inaudível:

— Respira, menino. Escuta o mundo antes de lutar contra ele.

Quando percebi, estava parado no meio da rua, exposto, distraído. O sol batia nas janelas, e as sombras pareciam mais longas. Eu havia esquecido que acabara de fugir de algo maior do que entendia.

Então um clarão tomou conta do beco. Um grito estourou.

— Então finalmente eu te achei!

O som me atingiu primeiro. Depois, o calor.  
Larissa dobrava a esquina com as mãos em brasa, o fogo dançando nos dedos. Corria na minha direção com o olhar que eu conhecia — o mesmo de antes, só que agora tomado de raiva.

O corpo inteiro se enrijeceu.  
A voz dentro da minha cabeça sumiu, como se nunca tivesse existido.

“Agora não dá pra filosofar”, pensei.  
A mente buscou freneticamente o que fazer.  
Eu mal sabia volitar. Com medo, as ilusões não vinham. Mover objetos parecia uma lembrança distante.  
Podia correr — sempre corri — mas não mais.  
Hoje, não.

O coração batia como tambor.  
— Chega de fugir — murmurei pra mim mesmo.

Ela se aproximava, brincando com as labaredas, confiante e furiosa. A luz refletia nos olhos, e, por um instante, vi a mesma garota de antes, antes de tudo desandar.

“Bom, minha amiga...” pensei, enquanto o calor aumentava.

— Eu sei que você não quer que eu lute — disse em voz baixa, para a presença silenciosa dentro da minha cabeça. — Mas, por favor, me diga alguma coisa. Porque, por mim mesmo, eu não faço ideia do que vou fazer.